% Figure: voltage profile along a radial distribution feeder. Under heavy load
% (blue) the voltage sags with distance as current draws a drop across the line
% resistance and reactance, approaching the lower statutory limit at the feeder
% end. With distributed rooftop PV exporting near the end of the feeder (orange)
% the flow reverses in the sunny midday and the voltage rises above nominal,
% approaching the upper limit, the reverse-flow overvoltage that limits how much
% PV a feeder can host. Values illustrative.
\begin{figure}[htbp]
\centering
\begin{tikzpicture}
\begin{axis}[
  width=13cm, height=8cm,
  xlabel={distance along feeder (fraction of length)},
  ylabel={voltage (per unit)},
  xmin=0, xmax=1, ymin=0.92, ymax=1.08,
  xtick={0,0.2,0.4,0.6,0.8,1.0},
  ytick={0.92,0.95,0.98,1.00,1.02,1.05,1.08},
  axis lines=left,
  tick label style={font=\scriptsize},
  label style={font=\small},
  legend style={font=\scriptsize, at={(0.03,0.03)}, anchor=south west, draw=none},
]
% statutory band +/- 5%
\addplot[enggray, thick, dashed] coordinates {(0,1.05)(1,1.05)};
\addplot[enggray, thick, dashed] coordinates {(0,0.95)(1,0.95)};
\node[anchor=west, font=\scriptsize, color=enggray] at (axis cs:0.02,1.062) {upper limit $+5\%$};
\node[anchor=west, font=\scriptsize, color=enggray] at (axis cs:0.02,0.938) {lower limit $-5\%$};
% nominal 1.0
\addplot[black, thin] coordinates {(0,1.00)(1,1.00)};
% heavy-load sag: voltage falls roughly linearly (cumulative I R drop) toward
% the end of the feeder
\addplot[engblue, very thick] coordinates {
  (0,1.00)(0.2,0.988)(0.4,0.976)(0.6,0.965)(0.8,0.956)(1.0,0.950)};
\addlegendentry{heavy evening load (voltage sags)}
% midday PV export: reverse flow lifts the far-end voltage above nominal
\addplot[engorange, very thick] coordinates {
  (0,1.00)(0.2,1.012)(0.4,1.024)(0.6,1.035)(0.8,1.044)(1.0,1.050)};
\addlegendentry{midday PV export (voltage rises)}
\end{axis}
\end{tikzpicture}
\caption[Voltage profile along a radial distribution feeder]{Voltage along a
radial distribution feeder under two opposite conditions. Under a heavy evening
load (blue) the current drawn along the feeder produces a cumulative drop
across the line resistance and reactance, so the voltage sags with distance and
is lowest at the far end, where it approaches the statutory lower limit. When
distributed rooftop solar exports near the end of the feeder in the sunny
midday (orange) the power flow reverses, and the same line impedance now lifts
the voltage above nominal, highest at the far end, approaching the upper limit.
The two limits together bound how much load and how much distributed generation
a feeder can carry, and the reverse-flow overvoltage is the constraint that
caps a feeder's photovoltaic hosting capacity, managed by voltage regulators,
by inverter volt-var control, or by reinforcing the line. Quantities
illustrative.}
\label{fig:vol04ch14:feeder-voltage}
\end{figure}
